% written by make_report_numbers.py
\begin{table}[t]
\centering
\caption{Training stability. Gradient: the size (norm) of the gradient of the loss in one update, before it is cut to 1; the median and the largest value of a run. Cut: the share of updates whose gradient was larger than 1 and was scaled down. Gradient spike: a gradient more than 5 times the median of the 100 updates before it. Loss spike: a loss more than 3 times the median of the last 50 updates on the same kind of batch. Where the runs differ, the lowest and the highest are given.}
\label{tab:stability}
\footnotesize\setlength{\tabcolsep}{4pt}
\begin{tabularx}{\linewidth}{@{}>{\raggedright\arraybackslash}Xlrcccccc@{}}
\toprule
 & & & & \multicolumn{2}{c}{Gradient} & & \multicolumn{2}{c}{Spikes} \\
\cmidrule(lr){5-6}\cmidrule(l){8-9}
Finder & Start & Updates & Runs & median & largest & Cut & gradient & loss \\
\midrule
EGNN & denoiser & 2,000 & 3 & 2.0--2.5 & 38--86 & 95--96\% & 9--32 & 1--13 \\
EGNN, no class weights & denoiser & 2,000 & 2 & 0.31--0.34 & 6.2--6.7 & 1--3\% & 3--4 & 5--11 \\
EGNN & random & 2,000 & 3 & 1.4--1.5 & 4.1--5.2 & 92--93\% & 0 & 0 \\
\addlinespace
EGNN & random & 20,000 & 3 & 1.5--1.6 & 55--82 & 85--88\% & 71--131 & 14--34 \\
Transformer & random & 20,000 & 3 & 0.9--1.0 & 4.9--5.5 & 40--50\% & 0 & 0--1 \\
CNN & random & 20,000 & 3 & 1.6--1.8 & 6.3--8.1 & 85\% & 0 & 0 \\
\bottomrule
\end{tabularx}
\end{table}
